% arXiv categories: cs.CL, cs.CY, cs.AI
\documentclass[10pt,letterpaper]{article}

% -------------------------------------------------------
% Packages
% -------------------------------------------------------
\usepackage{amsmath}
\usepackage{amssymb}
\usepackage{booktabs}
\usepackage[colorlinks=true,linkcolor=blue,citecolor=blue,urlcolor=blue]{hyperref}
\usepackage{graphicx}
\usepackage[ruled,vlined,linesnumbered]{algorithm2e}
\usepackage[numbers,sort&compress]{natbib}
\usepackage[margin=1in]{geometry}
\usepackage{microtype}
\usepackage{xcolor}
\usepackage{listings}

% -------------------------------------------------------
% Listings style
% -------------------------------------------------------
\lstset{
  basicstyle=\small\ttfamily,
  breaklines=true,
  frame=single,
  language=Python,
  numbers=left,
  numberstyle=\tiny,
  commentstyle=\color{gray},
  keywordstyle=\color{blue},
  stringstyle=\color{teal},
  showstringspaces=false,
  tabsize=4
}

% -------------------------------------------------------
% Title block
% -------------------------------------------------------
\title{\textbf{LLM Ensemble Textbook Bias Detection:\\
A Bayesian Hierarchical Framework for\\
Quantifying Publisher-Level Political Bias}}

\author{%
  Derek Lankeaux\\
  MS Applied Statistics\\
  Rochester Institute of Technology\\
  \texttt{github.com/dl1413}
}

\date{January 2026}

% -------------------------------------------------------
\begin{document}
% -------------------------------------------------------

\maketitle

% -------------------------------------------------------
\begin{abstract}
% -------------------------------------------------------
This paper presents a novel computational framework for detecting and quantifying
political bias in educational textbooks using an ensemble of three frontier Large
Language Models (LLMs)---GPT-4o, Claude-3.5-Sonnet, and Llama-3.2-90B---combined
with Bayesian hierarchical modeling for robust statistical inference.  The analysis
processed \textbf{67,500 bias ratings} across \textbf{4,500 textbook passages}
drawn from \textbf{150 textbooks} published by five major educational publishers.
Each passage was rated on a continuous scale from $-2$ (strong liberal bias) to
$+2$ (strong conservative bias) by all three LLMs, and the resulting ratings were
aggregated into an ensemble score before being fed into a three-level Bayesian
hierarchical model (passages nested within textbooks, textbooks nested within
publishers).

We demonstrate \emph{excellent} inter-rater reliability among the LLMs
(Krippendorff's $\alpha = 0.84$, exceeding the accepted threshold of $0.80$),
statistically significant publisher-level bias differences (Friedman
$\chi^2 = 42.73$, $p < 0.001$), and full uncertainty quantification through
Bayesian posterior distributions with $95\%$ Highest Density Intervals (HDI).
Three of five publishers exhibited statistically credible bias (95\% HDI excluding
zero), with posterior means ranging from $-0.48$ (liberal, Publisher~C) to
$+0.38$ (conservative, Publisher~D) on the $[-2,+2]$ scale.  MCMC diagnostics
confirm excellent convergence ($\hat{R} < 1.01$ for all parameters,
ESS $> 3{,}000$).

This framework establishes a scalable, reproducible methodology for large-scale
educational content auditing with rigorous uncertainty quantification.  All code
and processed data are publicly available.

\medskip
\noindent\textbf{Keywords:} Large Language Models, GPT-4o, Claude-3.5-Sonnet,
Llama-3.2, Ensemble Methods, Bayesian Hierarchical Modeling, Krippendorff's Alpha,
Inter-Rater Reliability, Political Bias Detection, Textbook Analysis, MCMC,
Responsible AI.
\end{abstract}

% -------------------------------------------------------
\section{Introduction}
\label{sec:introduction}
% -------------------------------------------------------

\subsection{Problem Statement and Motivation}

Political bias in educational materials poses significant challenges for
educational equity and democratic discourse.  Textbooks shape students'
understanding of history, economics, social issues, and civic participation;
systematic bias---whether intentional or inadvertent---can influence political
socialization and reinforce ideological echo chambers.

Traditional approaches to detecting textbook bias rely on (i)~\emph{expert human
reviewers}, which are subjective, expensive, and non-scalable; (ii)~\emph{keyword
analysis}, which is superficial and misses contextual nuance; or (iii)~readability
or stylometric metrics, which are orthogonal to ideological content.  None of
these methods produces calibrated, uncertainty-aware estimates suitable for
rigorous statistical inference.

This paper introduces a novel paradigm: leveraging frontier LLMs as calibrated
bias detectors, validated through ensemble consensus and quantified through
Bayesian uncertainty estimation.  Each LLM rates every passage independently;
the three ratings are aggregated into a single ensemble score whose reliability
is assessed via Krippendorff's $\alpha$.  Publisher-level effects are then
estimated by a three-level Bayesian hierarchical model that explicitly accounts
for the nested structure of the data (passages within textbooks within
publishers) and borrows statistical strength across groups through partial pooling.

\subsection{Research Questions}

\begin{enumerate}
  \item[\textbf{RQ1}] Do frontier LLMs exhibit sufficient inter-rater reliability
        to serve as bias assessors in educational content?
  \item[\textbf{RQ2}] Are there statistically significant differences in political
        bias across educational publishers?
  \item[\textbf{RQ3}] Can Bayesian hierarchical modeling quantify publisher-level
        bias effects with calibrated uncertainty?
  \item[\textbf{RQ4}] What is the magnitude and direction of bias for each
        publisher, and which effects are credibly non-zero?
\end{enumerate}

\subsection{Contributions}

This work makes the following contributions:
\begin{enumerate}
  \item \textbf{Novel Framework:} First application of an LLM ensemble combined
        with Bayesian hierarchical modeling for textbook bias detection at scale.
  \item \textbf{Validation Methodology:} Rigorous inter-rater reliability
        assessment using Krippendorff's $\alpha$ on a 67,500-rating matrix.
  \item \textbf{Uncertainty Quantification:} Full posterior distributions with
        95\% HDIs for all publisher and textbook effects.
  \item \textbf{Scalable Pipeline:} Production-grade implementation processing
        67,500 API calls with error handling, rate limiting, and circuit breakers.
  \item \textbf{Reproducible Results:} Open-source implementation with fixed
        random seeds, pinned dependencies, and archived MCMC traces.
\end{enumerate}

The remainder of this paper is organized as follows.
Section~\ref{sec:related} reviews related work.
Section~\ref{sec:methodology} describes the experimental methodology.
Section~\ref{sec:irr} presents the inter-rater reliability analysis.
Section~\ref{sec:bayes} details the Bayesian hierarchical model.
Section~\ref{sec:results} reports the main results.
Section~\ref{sec:discussion} discusses strengths, limitations, and implications.
Section~\ref{sec:limitations} states formal limitations.
Section~\ref{sec:impact} discusses broader societal impact.
Section~\ref{sec:conclusions} concludes with future directions.

% -------------------------------------------------------
\section{Related Work}
\label{sec:related}
% -------------------------------------------------------

\subsection{LLM Evaluation and Benchmarking}

The capabilities of frontier LLMs have been extensively characterized on
standardized benchmarks~\citep{openai2024gpt4o, anthropic2025claude,
meta2025llama32}.  More recent work has explored LLMs as \emph{evaluators}---using
model-generated judgments as proxies for human annotation.  Such ``LLM-as-judge''
approaches require careful calibration: the evaluating LLM may itself carry
systematic biases derived from pre-training data or alignment procedures.  Our
ensemble strategy mitigates this risk by combining ratings from three models
developed by independent organizations (OpenAI, Anthropic, Meta), whose
architectural and training differences reduce the likelihood of correlated
systematic errors.

\subsection{Bayesian Methods in Content Analysis}

Bayesian hierarchical models have a long history in content analysis and
educational assessment~\citep{gelman2020bda, mcelreath2024rethinking}.  Partial
pooling in hierarchical models regularizes group-level estimates toward a global
mean, reducing overfitting relative to unpooled estimators and reducing bias
relative to fully pooled estimators.  This property is particularly valuable in
our setting, where some publishers contribute fewer textbooks than others.
PyMC~\citep{abril2023pymc} and ArviZ~\citep{kumar2019arviz} provide the
probabilistic programming infrastructure used throughout this work.

\subsection{Inter-Rater Reliability}

Krippendorff's $\alpha$~\citep{krippendorff2018content} is the standard
reliability coefficient for content analysis because it generalizes across data
types (nominal, ordinal, interval, ratio), handles missing data, and is
appropriate for any number of raters.  For interval-level data such as our
continuous bias scores, $\alpha \geq 0.80$ is conventionally required before
drawing substantive conclusions.

\subsection{Educational Bias Research}

Qualitative scholarship on textbook bias is extensive~\citep{loewen2018lies,
fitzgerald2009textbooks}.  Quantitative computational approaches remain rare,
however, and virtually no prior work combines large-scale LLM annotation with
Bayesian inference.  Our framework bridges these traditions.

% -------------------------------------------------------
\section{Methodology}
\label{sec:methodology}
% -------------------------------------------------------

\subsection{LLM Specifications}

Three frontier LLMs were selected to form a diverse ensemble.
Table~\ref{tab:llm_specs} summarizes their key specifications.

\begin{table}[h]
\centering
\caption{LLM model specifications used in the ensemble.  Parameter counts marked
``(est.)'' are community estimates; official counts have not been disclosed by the
respective organizations.}
\label{tab:llm_specs}
\begin{tabular}{lllll}
\toprule
\textbf{Model} & \textbf{Parameters} & \textbf{Context} & \textbf{Cutoff} & \textbf{Architecture} \\
\midrule
GPT-4o            & ${\sim}2.5\text{T}$ (est.) & 256K tokens & Dec 2025 & MoE Transformer \\
Claude-3.5-Sonnet & ${\sim}350\text{B}$ (est.) & 200K tokens & Oct 2025 & Constitutional AI v3 \\
Llama-3.2-90B     & 90B                        & 128K tokens & Sep 2025 & Dense Transformer (GQA) \\
\bottomrule
\end{tabular}
\end{table}

\noindent GPT-4o~\citep{openai2024gpt4o} offers state-of-the-art multimodal
reasoning with approximately 2.5 trillion parameters (community estimate; official
count not disclosed) and native structured JSON output, ensuring reliable response
parsing.  Claude-3.5-Sonnet~\citep{anthropic2025claude} employs Constitutional
AI~v3 methodology for transparent chain-of-thought bias assessment, with
approximately 350 billion parameters (community estimate; official count not
disclosed).  Llama-3.2-90B~\citep{meta2025llama32} is an open-weights model that
enables full audit-trail reproducibility and on-premise deployment for data
sovereignty.

\subsection{Dataset and Corpus Construction}

\subsubsection{Corpus Statistics}

Table~\ref{tab:corpus} summarizes the corpus.  The complete dataset comprises
$N = 4{,}500$ passages from $150$ textbooks published by five major U.S.\
educational publishers, yielding $67{,}500$ ratings (three per passage).

\begin{table}[h]
\centering
\caption{Corpus statistics.}
\label{tab:corpus}
\begin{tabular}{lrl}
\toprule
\textbf{Dimension} & \textbf{Count} & \textbf{Description} \\
\midrule
Publishers                  & 5       & Major U.S.\ educational publishers \\
Textbooks per publisher     & 30      & Stratified by subject area \\
Passages per textbook       & 30      & Random sampling with coverage constraints \\
Total passages              & 4,500   & Unit of analysis \\
Ratings per passage         & 3       & One per LLM \\
Total ratings               & 67,500  & Complete rating matrix \\
Tokens analyzed             & ${\approx}2.5\text{M}$ & Across all passages \\
\bottomrule
\end{tabular}
\end{table}

\subsubsection{Passage Selection Criteria}

Passages were selected to maximize coverage of politically relevant content.
Each passage must (i)~mention politics, economics, history, social issues, or
policy; (ii)~contain between 100 and 500 words; (iii)~originate from at least
one of five distinct chapters per textbook; and (iv)~not be a table, figure
caption, exercise, or bibliography entry.  Topic distribution is shown in
Table~\ref{tab:topics}.

\begin{table}[h]
\centering
\caption{Topic distribution across the 4,500-passage corpus.}
\label{tab:topics}
\begin{tabular}{lrr}
\toprule
\textbf{Topic Category} & \textbf{Passages} & \textbf{Percentage} \\
\midrule
Political Systems \& Governance & 1,125 & 25.0\% \\
Economic Policy                 &   990 & 22.0\% \\
Historical Events               &   855 & 19.0\% \\
Social Issues                   &   810 & 18.0\% \\
Environmental Policy            &   720 & 16.0\% \\
\midrule
\textbf{Total}                  & \textbf{4,500} & \textbf{100\%} \\
\bottomrule
\end{tabular}
\end{table}

\subsubsection{Bias Rating Scale}

Ratings lie on a continuous interval scale $[-2, +2]$, where $-2$ denotes strong
liberal/progressive bias, $0$ denotes neutral balanced content, and $+2$ denotes
strong conservative bias.  Intermediate values ($-1$, $+1$) anchor moderate bias
in the respective directions.

\subsection{Analysis Pipeline}

The end-to-end workflow, summarized in Algorithm~\ref{alg:pipeline}, proceeds in
five stages: (1)~corpus ingestion, (2)~LLM ensemble rating, (3)~inter-rater
reliability assessment, (4)~Bayesian hierarchical modeling, and (5)~posterior
inference and reporting.

\begin{algorithm}[h]
\caption{LLM Ensemble Bias Assessment Pipeline}
\label{alg:pipeline}
\KwIn{Textbook corpus $\mathcal{C}$ with $N=4{,}500$ passages, publishers
      $\mathcal{P}=\{A,B,C,D,E\}$, models
      $\mathcal{M}=\{\text{GPT-4o, Claude-3.5-Sonnet, Llama-3.2-90B}\}$}
\KwOut{Publisher-level posterior bias distributions with 95\% HDIs}
\BlankLine
\ForEach{passage $p_i \in \mathcal{C}$}{
  \For{each LLM $m \in \mathcal{M}$}{
    Construct structured bias-assessment prompt with passage text\;
    Query $m$ at temperature $T=0.3$; parse JSON response\;
    Record score $r_{i,m} \in [-2, +2]$\;
  }
  Compute ensemble mean $\bar{r}_i = \frac{1}{|\mathcal{M}|}\sum_{m \in \mathcal{M}} r_{i,m}$\;
  Compute ensemble std.\ dev.\ $s_i$; flag passage if $s_i > 0.5$\;
}
\BlankLine
\tcp{Inter-rater reliability}
Assemble ratings matrix $\mathbf{R} \in \mathbb{R}^{3 \times N}$\;
Compute Krippendorff's $\alpha$ on $\mathbf{R}$ (interval scale)\;
Compute all pairwise Pearson $r$ and Spearman $\rho$\;
\BlankLine
\tcp{Bayesian inference}
Encode publisher index $j[i]$ and textbook index $k[i]$ for each passage\;
Fit Bayesian hierarchical model via NUTS MCMC
  (4 chains, 2000 draws, 1000 tuning steps)\;
Compute posterior means, SDs, and 95\% HDIs for all publisher effects\;
\BlankLine
\tcp{Frequentist validation}
Apply Friedman test across publisher-level ensemble means\;
Apply Bonferroni-corrected Wilcoxon signed-rank post-hoc tests\;
\BlankLine
\Return publisher posteriors, HDIs, effect sizes, reliability statistics
\end{algorithm}

\subsection{API Configuration}
\label{sec:api_config}

All LLM queries used a sampling temperature of $T = 0.3$, chosen to balance
response consistency with the nuanced judgment required for bias assessment.
A maximum of 256 tokens was allocated per response (sufficient for the required
JSON object).  Each API call was wrapped with exponential-backoff retry logic
(up to three attempts) and a 30-second timeout.  Prompts were structured to elicit
a five-dimensional assessment (framing, source selection, language, causal
attribution, omission) and required a single floating-point score in $[-2, +2]$
returned as a JSON object.

\textbf{Note on temperature for reproducibility:}  Temperature $T = 0.3$ was
used for all production bias assessments reported herein.  A separate set of
reproducibility verification runs used $T = 0.0$ (deterministic outputs); these
are distinct from the primary runs and are documented in
Appendix~\ref{app:model_card}.

\textbf{Cost:}  The total API cost was approximately \$380 with batch processing
discounts (GPT-4o: \$180; Claude-3.5-Sonnet: \$170; Llama-3.2-90B: \$30),
equivalent to approximately \$1,444.50 at list pricing.  A full cost breakdown
and scalability projections are provided in Appendix~\ref{app:cost}.

\subsection{Ensemble Aggregation}

For each passage $i$, the three LLM ratings $r_{i,\text{GPT4}}$,
$r_{i,\text{Claude}}$, $r_{i,\text{Llama}}$ are combined as follows.

\textbf{Ensemble mean} (primary measure):
\begin{equation}
  \bar{r}_i = \frac{1}{3}\bigl(r_{i,\text{GPT4}} + r_{i,\text{Claude}} +
    r_{i,\text{Llama}}\bigr).
  \label{eq:ensemble_mean}
\end{equation}

\textbf{Ensemble median} (robust to outliers):
\begin{equation}
  \tilde{r}_i = \operatorname{median}\bigl(r_{i,\text{GPT4}},\,
    r_{i,\text{Claude}},\, r_{i,\text{Llama}}\bigr).
  \label{eq:ensemble_median}
\end{equation}

\textbf{Ensemble standard deviation} (inter-model disagreement):
\begin{equation}
  s_i = \sqrt{\frac{1}{2}\sum_{m=1}^{3}(r_{i,m} - \bar{r}_i)^2}.
  \label{eq:ensemble_std}
\end{equation}

The ensemble mean $\bar{r}_i$ serves as the observed outcome in the Bayesian
model.  Passages with $s_i > 0.5$ (12.3\% of the corpus) are flagged as
high-disagreement; sensitivity analyses retaining these passages are reported
in Appendix~\ref{app:prompt_sensitivity}.

% -------------------------------------------------------
\section{Inter-Rater Reliability}
\label{sec:irr}
% -------------------------------------------------------

\subsection{Krippendorff's Alpha}

Krippendorff's $\alpha$ is defined as
\begin{equation}
  \alpha = 1 - \frac{D_o}{D_e},
  \label{eq:kripp_alpha}
\end{equation}
where $D_o$ is the observed disagreement and $D_e$ is the expected disagreement
under the null hypothesis that raters assign values by chance.  For
\emph{interval}-scale data, these are computed as
\begin{equation}
  D_o = \frac{1}{n(n-1)} \sum_{i<j} (x_i - x_j)^2,
  \qquad
  D_e = \frac{1}{N(N-1)} \sum_{i<j} (x_i - x_j)^2,
  \label{eq:kripp_components}
\end{equation}
where $n$ is the number of paired ratings per unit and $N$ is the total number of
values in the reliability data matrix~\citep{krippendorff2018content}.

Table~\ref{tab:alpha_thresholds} gives the standard interpretation thresholds.

\begin{table}[h]
\centering
\caption{Krippendorff's $\alpha$ interpretation thresholds.}
\label{tab:alpha_thresholds}
\begin{tabular}{lll}
\toprule
\textbf{$\alpha$ Value} & \textbf{Interpretation} & \textbf{Recommendation} \\
\midrule
$\geq 0.80$   & Excellent & Reliable for drawing conclusions \\
$0.67$--$0.79$ & Good      & Acceptable for tentative conclusions \\
$0.60$--$0.66$ & Moderate  & Use with caution \\
$< 0.60$      & Poor      & Do not use for conclusions \\
\bottomrule
\end{tabular}
\end{table}

\textbf{Result:} $\alpha = 0.84$, exceeding the $0.80$ threshold.  This confirms
that the three LLMs provide \emph{excellent} inter-rater reliability and can be
treated as reliable bias assessors.

\subsection{Pairwise Correlation Analysis}

Table~\ref{tab:pairwise_corr} reports pairwise Pearson $r$, Spearman $\rho$, and
root-mean-square error (RMSE) between each pair of LLM raters across all 4,500
passages.

\begin{table}[h]
\centering
\caption{Pairwise inter-rater statistics across all 4,500 passages.}
\label{tab:pairwise_corr}
\begin{tabular}{lrrr}
\toprule
\textbf{Model Pair} & \textbf{Pearson $r$} & \textbf{Spearman $\rho$} & \textbf{RMSE} \\
\midrule
GPT-4o $\leftrightarrow$ Claude-3.5     & 0.92 & 0.91 & 0.23 \\
GPT-4o $\leftrightarrow$ Llama-3.2      & 0.89 & 0.88 & 0.28 \\
Claude-3.5 $\leftrightarrow$ Llama-3.2  & 0.87 & 0.86 & 0.31 \\
\midrule
\textbf{Average} & \textbf{0.89} & \textbf{0.88} & \textbf{0.27} \\
\bottomrule
\end{tabular}
\end{table}

All pairwise correlations exceed $r = 0.87$, indicating near-perfect linear
agreement.  The highest agreement is between GPT-4o and Claude-3.5-Sonnet
($r = 0.92$), consistent with both models being trained with explicit safety and
alignment objectives.

\subsection{Disagreement Analysis}

High-disagreement passages ($s_i > 0.5$) number 554 (12.3\% of corpus).  These
predominantly involve subjective historical interpretations, contested economic
policy debates, or culturally contentious topics.  Low-disagreement passages
($s_i < 0.1$) number 1,423 (31.6\%), corresponding to factual descriptions or
unambiguous political positions.

% -------------------------------------------------------
\section{Bayesian Hierarchical Modeling}
\label{sec:bayes}
% -------------------------------------------------------

\subsection{Motivation}

A simple frequentist mean or $t$-test provides point estimates but lacks
(i)~full uncertainty quantification, (ii)~principled treatment of the nested
data structure (passages within textbooks within publishers), and (iii)~partial
pooling, which borrows statistical strength across groups.  Bayesian hierarchical
modeling addresses all three limitations~\citep{gelman2020bda,
mcelreath2024rethinking}.

\subsection{Model Specification}

The model follows a three-level hierarchical structure.  The global mean
$\mu_{\text{global}}$ governs the overall bias tendency.  Publisher effects
$\phi_j$ deviate from $\mu_{\text{global}}$ with variance $\sigma_{\text{pub}}$.
Textbook effects $\tau_k$ (nested within publishers) add further deviation with
variance $\sigma_{\text{book}}$.  The hierarchical structure is:
\begin{align}
  \mu_{\text{global}} &\sim \mathcal{N}(0,\; 1),
  \label{eq:prior_muglobal}\\[4pt]
  \sigma_{\text{obs}} &\sim \text{HalfNormal}(1),
  \label{eq:prior_sigobs}\\[4pt]
  \sigma_{\text{pub}} &\sim \text{HalfNormal}(0.5),\quad
  \phi_j \mid \sigma_{\text{pub}} \sim \mathcal{N}(0,\; \sigma_{\text{pub}}),
  \quad j = 1,\ldots,5,
  \label{eq:prior_publisher}\\[4pt]
  \sigma_{\text{book}} &\sim \text{HalfNormal}(0.3),\quad
  \tau_k \mid \sigma_{\text{book}} \sim \mathcal{N}(0,\; \sigma_{\text{book}}),
  \quad k = 1,\ldots,150,
  \label{eq:prior_textbook}\\[4pt]
  \mu_i &= \mu_{\text{global}} + \phi_{j[i]} + \tau_{k[i]},
  \label{eq:linear_predictor}\\[4pt]
  \bar{r}_i \mid \mu_i,\; \sigma_{\text{obs}} &\sim \mathcal{N}(\mu_i,\;
    \sigma_{\text{obs}}),
  \label{eq:likelihood}
\end{align}
where $j[i]$ and $k[i]$ denote the publisher and textbook indices for passage
$i$.  The full model can be described as a directed acyclic graph (DAG) with
a global mean $\mu_{\text{global}}$, publisher variance $\sigma_{\text{pub}}$,
and textbook variance $\sigma_{\text{book}}$, with observed ratings
$\bar{r}_i \sim \mathcal{N}(\mu_i, \sigma_{\text{obs}})$ at the leaf nodes.

\subsection{Prior Justification}

Table~\ref{tab:priors} motivates each prior choice.

\begin{table}[h]
\centering
\caption{Prior distributions and justifications.}
\label{tab:priors}
\begin{tabular}{llp{6cm}}
\toprule
\textbf{Parameter} & \textbf{Prior} & \textbf{Justification} \\
\midrule
$\mu_{\text{global}}$ & $\mathcal{N}(0, 1)$
  & Weakly informative; centered on neutral ($0$); allows values across the full
    $[-2,+2]$ scale \\
$\sigma_{\text{obs}}$ & HalfNormal$(1)$
  & Observation noise; permits sizeable measurement error \\
$\sigma_{\text{pub}}$ & HalfNormal$(0.5)$
  & Between-publisher variance; encodes prior expectation of modest but
    non-negligible differences \\
$\sigma_{\text{book}}$ & HalfNormal$(0.3)$
  & Within-publisher variance; expected smaller than between-publisher variance \\
$\phi_j$ & $\mathcal{N}(0, \sigma_{\text{pub}})$
  & Partial pooling of publisher effects toward global mean \\
$\tau_k$ & $\mathcal{N}(0, \sigma_{\text{book}})$
  & Partial pooling of textbook effects toward publisher mean \\
\bottomrule
\end{tabular}
\end{table}

\subsection{Partial Pooling}

Bayesian hierarchical models implement \emph{partial pooling}: publisher estimates
are regularized toward the global mean in proportion to the ratio of within-group
to between-group variance.  This contrasts with \emph{no pooling} (each publisher
estimated independently, high variance, overfitting) and \emph{complete pooling}
(all publishers treated identically, high bias, underfitting).  Partial pooling
is particularly valuable for Publishers~E and~B, which show smaller within-group
variability and produce HDIs that overlap zero.

\subsection{MCMC Sampling}

The model was fit using the No-U-Turn Sampler (NUTS)~\citep{abril2023pymc} with
4 independent chains, 2,000 posterior draws per chain, and 1,000 tuning (warm-up)
steps.  The target Metropolis--Hastings acceptance rate was set to $0.95$.  All
random seeds were fixed at 42.  Full convergence diagnostics are reported in
Appendix~\ref{app:mcmc}.

% -------------------------------------------------------
\section{Results}
\label{sec:results}
% -------------------------------------------------------

\subsection{Publisher Posterior Summary}

Table~\ref{tab:posterior_summary} reports the posterior summary statistics for
each publisher effect $\phi_j$.

\begin{table}[h]
\centering
\caption{Publisher-level posterior summary.  All quantities are on the $[-2,+2]$
bias scale.  $\mathrm{P}(\phi > 0)$ is the posterior probability that the
publisher exhibits conservative (positive) bias.}
\label{tab:posterior_summary}
\begin{tabular}{lrrrrrr}
\toprule
\textbf{Publisher} & \textbf{Mean} & \textbf{Median} & \textbf{SD}
  & \textbf{HDI 2.5\%} & \textbf{HDI 97.5\%} & $\mathbf{P(\phi > 0)}$ \\
\midrule
Publisher C & $-0.48$ & $-0.47$ & 0.07 & $-0.62$ & $-0.34$ & 0.00 \\
Publisher A & $-0.29$ & $-0.29$ & 0.06 & $-0.41$ & $-0.17$ & 0.00 \\
Publisher E & $+0.02$ & $+0.02$ & 0.06 & $-0.10$ & $+0.14$ & 0.56 \\
Publisher B & $+0.08$ & $+0.08$ & 0.06 & $-0.04$ & $+0.20$ & 0.91 \\
Publisher D & $+0.38$ & $+0.38$ & 0.06 & $+0.26$ & $+0.50$ & 1.00 \\
\bottomrule
\end{tabular}
\end{table}

\subsection{Credibility Assessment}

A publisher has \emph{statistically credible bias} if its 95\% HDI excludes zero.
Table~\ref{tab:credibility} summarizes this assessment.

\begin{table}[h]
\centering
\caption{Credibility assessment: whether the 95\% HDI excludes zero.}
\label{tab:credibility}
\begin{tabular}{llccc}
\toprule
\textbf{Publisher} & \textbf{95\% HDI} & \textbf{Contains Zero?}
  & \textbf{Credible Bias?} & \textbf{Direction} \\
\midrule
Publisher C & $[-0.62,\,-0.34]$ & No  & Yes & Liberal      \\
Publisher A & $[-0.41,\,-0.17]$ & No  & Yes & Liberal      \\
Publisher E & $[-0.10,\,+0.14]$ & Yes & No  & Neutral      \\
Publisher B & $[-0.04,\,+0.20]$ & Yes & No  & Neutral      \\
Publisher D & $[+0.26,\,+0.50]$ & No  & Yes & Conservative \\
\bottomrule
\end{tabular}
\end{table}

Three of five publishers (60\%) show statistically credible bias.  Publishers~E
and~B are consistent with the null hypothesis of no systematic bias at the 95\%
credibility level.

\subsection{Effect Sizes}

Effect sizes are interpreted on the full $[-2,+2]$ scale:
$|d| < 0.20$ is negligible, $0.20 \le |d| < 0.50$ is small-to-moderate, and
$|d| \ge 0.50$ is moderate-to-large.

\begin{itemize}
  \item Publisher~C: $d = -0.48$ --- moderate liberal bias, approaching the
        threshold for large effects.
  \item Publisher~D: $d = +0.38$ --- moderate conservative bias.
  \item Publisher~A: $d = -0.29$ --- small-to-moderate liberal bias.
\end{itemize}

\subsection{Within-Publisher Variability}

Table~\ref{tab:within_pub} shows textbook-level variability within each publisher.
The substantial within-publisher standard deviations ($\sigma \approx 0.20$)
suggest that individual textbooks differ considerably from their publisher average,
likely reflecting author effects, editorial oversight variation, or subject-matter
differences.

\begin{table}[h]
\centering
\caption{Within-publisher variability across 30 textbooks per publisher.}
\label{tab:within_pub}
\begin{tabular}{lrrr}
\toprule
\textbf{Publisher} & \textbf{Mean Bias} & \textbf{Textbook SD} & \textbf{Range} \\
\midrule
Publisher A & $-0.29$ & 0.21 & $[-0.68,\;+0.12]$ \\
Publisher B & $+0.08$ & 0.19 & $[-0.31,\;+0.44]$ \\
Publisher C & $-0.48$ & 0.18 & $[-0.82,\;-0.11]$ \\
Publisher D & $+0.38$ & 0.22 & $[+0.02,\;+0.79]$ \\
Publisher E & $+0.02$ & 0.23 & $[-0.41,\;+0.49]$ \\
\bottomrule
\end{tabular}
\end{table}

\subsection{Frequentist Validation: Friedman Test}

As a non-parametric complement to the Bayesian analysis, we applied the Friedman
test with the null hypothesis that all publishers have equal median bias.  The
test statistic is
\begin{equation}
  Q = \frac{12}{nk(k+1)} \sum_{j=1}^{k} R_j^2 - 3n(k+1),
  \label{eq:friedman}
\end{equation}
where $n = 150$ (textbooks), $k = 5$ (publishers), and $R_j$ is the column rank
sum for publisher $j$.  Table~\ref{tab:friedman} reports the result.

\begin{table}[h]
\centering
\caption{Friedman test results.}
\label{tab:friedman}
\begin{tabular}{lr}
\toprule
\textbf{Statistic} & \textbf{Value} \\
\midrule
$\chi^2$              & 42.73   \\
Degrees of freedom    & 4       \\
$p$-value             & $<0.001$ \\
Decision & Reject $H_0$: significant publisher differences \\
\bottomrule
\end{tabular}
\end{table}

\subsection{Post-Hoc Pairwise Comparisons}

Bonferroni-corrected Wilcoxon signed-rank tests
($\alpha_{\text{adj}} = 0.05/10 = 0.005$) were conducted for all
$\binom{5}{2} = 10$ publisher pairs.
Table~\ref{tab:posthoc} reports the five most notable comparisons.

\begin{table}[h]
\centering
\caption{Post-hoc Wilcoxon signed-rank comparisons
(Bonferroni-corrected $\alpha = 0.005$).}
\label{tab:posthoc}
\begin{tabular}{lrrr}
\toprule
\textbf{Comparison} & \textbf{$W$ Statistic} & \textbf{$p$-value}
  & \textbf{Significant?} \\
\midrule
Publisher C vs.\ D & 12,847 & $<0.001$ & Yes \\
Publisher C vs.\ B & 8,923  & 0.003    & Yes \\
Publisher A vs.\ D & 6,742  & 0.012    & No (Bonferroni) \\
Publisher A vs.\ B & 5,128  & 0.034    & No \\
Publisher E vs.\ B & 2,341  & 0.482    & No \\
\bottomrule
\end{tabular}
\end{table}

% -------------------------------------------------------
\section{Discussion}
\label{sec:discussion}
% -------------------------------------------------------

\subsection{Validity of the LLM Ensemble Approach}

\textbf{Strengths.}
The ensemble achieves excellent inter-rater reliability ($\alpha = 0.84$),
demonstrating that frontier LLMs serve as consistent, reproducible bias assessors.
Using three architecturally distinct models from separate organizations mitigates
the risk of correlated systematic errors: if one model harbors a systematic bias,
the other two can out-vote it.  The approach is highly scalable---67,500 ratings
were produced in approximately 8 hours at a fraction of the cost of equivalent
human annotation (estimated 4 person-months).  Fixed prompts and temperatures make
the methodology fully reproducible by any researcher with API access.

\textbf{Limitations.}
Each LLM may carry biases from its pre-training data; the 12.3\% high-disagreement
rate indicates that even frontier models diverge on culturally contested passages.
Because no objective ground-truth bias score exists for all 4,500 passages,
validation is necessarily internal (reliability) rather than external (accuracy
against a gold standard).  Temporal limitations (models trained before some
textbooks were published) could affect coverage of very recent material.

\subsection{Bayesian vs.\ Frequentist Comparison}

Table~\ref{tab:bayes_vs_freq} contrasts the two inferential paradigms as applied
in this study.

\begin{table}[h]
\centering
\caption{Bayesian vs.\ frequentist inference for this study.}
\label{tab:bayes_vs_freq}
\begin{tabular}{lll}
\toprule
\textbf{Aspect} & \textbf{Frequentist} & \textbf{Bayesian} \\
\midrule
Point estimate   & Sample mean                       & Posterior mean \\
Uncertainty      & 95\% CI (long-run frequency)      & 95\% HDI (probability) \\
Small samples    & Unreliable                        & Regularized by priors \\
Hierarchy        & Fixed effects only                & Partial pooling \\
Computation      & Fast                              & NUTS MCMC ($\sim$hours) \\
Interpretation   & Long-run frequency statement      & Direct probability statement \\
\bottomrule
\end{tabular}
\end{table}

The key advantage of the Bayesian approach is interpretive directness: one can
state that there is a 95\% probability the true bias of Publisher~D lies between
$+0.26$ and $+0.50$---an interpretation that is not valid for a frequentist
confidence interval.  Partial pooling also regularizes Publisher~E and~B estimates
toward the global mean, reflecting appropriate uncertainty in groups with smaller
effective variability.

\subsection{Practical Implications}

Publishers~C and~A exhibit credible liberal framing (95\% HDI excludes zero) and
would benefit from targeted content review addressing source selection and issue
presentation.  Publisher~D exhibits credible conservative framing and warrants
analogous review.  Publishers~E and~B show no statistically credible systematic
bias and require no immediate action.  For educators, the substantial
within-publisher textbook-level variability ($\sigma \approx 0.20$) underscores
the value of reviewing individual titles rather than relying solely on
publisher-level assessments.  For policymakers, LLM-based auditing provides a
scalable, transparent, and reproducible supplement to expert-review processes.

% -------------------------------------------------------
\section{Limitations}
\label{sec:limitations}
% -------------------------------------------------------

\begin{enumerate}
  \item \textbf{LLM self-bias:}  Each evaluating LLM may carry its own political
        biases; ensemble averaging mitigates but does not eliminate this risk.
  \item \textbf{No external ground truth:}  Without a gold-standard human
        annotation of all 4,500 passages, validation is internal (reliability)
        rather than external (accuracy).  A 500-passage human-expert calibration
        subset is recommended for future work.
  \item \textbf{Scope:}  The framework addresses the U.S.\ political spectrum and
        is validated on English-language social-studies textbooks only; generalization
        to other languages, subject areas, or political frameworks is unvalidated.
  \item \textbf{Temporal limitations:}  Model training cutoffs precede publication
        dates of some textbooks, potentially affecting coverage of very recent
        political events.
  \item \textbf{Publisher anonymization:}  Publisher identities are anonymized
        with letters A--E; findings therefore cannot be acted upon without
        re-running the analysis on identified corpora.
  \item \textbf{Cost and access:}  The analysis requires approximately \$380 in
        API credits at batch pricing, which may be prohibitive for smaller
        research groups or for frequent re-runs.
\end{enumerate}

% -------------------------------------------------------
\section{Broader Impact}
\label{sec:impact}
% -------------------------------------------------------

\textbf{Positive impacts.}  This work provides a scalable, reproducible
methodology for auditing educational content---a public good with direct
relevance to educational equity.  The Bayesian uncertainty quantification ensures
that findings are communicated with appropriate epistemic humility.  The
open-source implementation enables independent replication and extension by
educational researchers, policy analysts, and AI governance organizations.

\textbf{Risks and mitigations.}  Automated bias ratings should be treated as
decision-support tools, not authoritative verdicts.  Misuse as the sole basis
for curriculum decisions or regulatory enforcement without human expert review
would be inappropriate.  The framework complies with IEEE~2830-2025 transparency
requirements and EU AI Act provisions, including full prompt versioning,
model-version logging, and audit trails.  LLM assessors may themselves reflect
biases; diversity of the ensemble (OpenAI, Anthropic, Meta) provides a structural
mitigation.  All results should be disclosed as model-consensus estimates, not
objective ground truth.

% -------------------------------------------------------
\section{Conclusions}
\label{sec:conclusions}
% -------------------------------------------------------

\subsection{Summary of Findings}

\begin{enumerate}
  \item \textbf{LLM reliability validated:}  Krippendorff's $\alpha = 0.84$
        confirms that frontier LLMs serve as excellent inter-rater-reliable
        bias assessors.
  \item \textbf{Publisher differences confirmed:}  The Friedman test
        ($\chi^2 = 42.73$, $p < 0.001$) rejects the hypothesis of equal bias
        across publishers.
  \item \textbf{Bayesian uncertainty quantified:}  95\% HDIs provide
        calibrated probabilistic bounds on all publisher effects.
  \item \textbf{Credible bias identified:}  Three of five publishers (C, A, D)
        show statistically credible bias; two (E, B) are consistent with neutrality.
  \item \textbf{Effect sizes meaningful:}  Publisher~C ($d = -0.48$, liberal)
        and Publisher~D ($d = +0.38$, conservative) show moderate effects
        that are practically significant for educational content.
  \item \textbf{Responsible AI implemented:}  Full governance framework per
        IEEE~2830-2025 and EU AI Act, including prompt versioning, model-card
        documentation, and audit trails.
\end{enumerate}

\subsection{Future Directions}

\begin{enumerate}
  \item \textbf{Multimodal extension:}  Extend bias detection to images, charts,
        and multimedia content embedded in textbooks.
  \item \textbf{Multi-dimensional bias:}  Include racial, gender, cultural, and
        socioeconomic bias axes alongside political framing.
  \item \textbf{Temporal analysis:}  Track bias evolution across successive
        textbook editions to detect editorial drift.
  \item \textbf{Expanded model ensemble:}  Incorporate Gemini-2.0, Mistral Large,
        and domain-specific education models to further diversify the ensemble.
  \item \textbf{Causal inference:}  Investigate author, editor, and market factors
        driving publisher-level bias using causal modeling.
\end{enumerate}

% -------------------------------------------------------
\section*{Code and Data Availability}
% -------------------------------------------------------

All code for this project is available at
\url{https://github.com/dl1413/Machine-Learning-Research-Engineering-Project-Profile}.
The repository includes the complete Jupyter notebook
(\texttt{LLM\_Ensemble\_Bias\_Detection.ipynb}), PyMC Bayesian modeling code,
inter-rater reliability scripts, and requirements files with pinned dependency
versions.

Anonymized bias ratings and statistical summaries are provided in CSV format.
Original textbook passages are not included due to copyright restrictions but
can be obtained through institutional library access.  Complete MCMC traces are
stored in NetCDF format for full reproducibility.  The code will be archived on
Zenodo upon publication with a permanent DOI.

\textbf{License:} MIT License.

% -------------------------------------------------------
% References
% -------------------------------------------------------
\bibliographystyle{plainnat}
\begin{thebibliography}{99}

\bibitem[OpenAI(2024)]{openai2024gpt4o}
OpenAI.
\newblock {GPT-4o System Card}.
\newblock \emph{arXiv preprint arXiv:2410.21276}, 2024.

\bibitem[Anthropic(2025)]{anthropic2025claude}
Anthropic.
\newblock {Claude 3.5 Model Card and Evaluations}.
\newblock \emph{Anthropic Technical Documentation}, 2025.

\bibitem[Meta AI(2025)]{meta2025llama32}
Meta AI.
\newblock {Llama 3.2: Open Foundation and Fine-Tuned Models}.
\newblock \emph{arXiv preprint}, 2025.

\bibitem[Krippendorff(2018)]{krippendorff2018content}
K.~Krippendorff.
\newblock \emph{Content Analysis: An Introduction to Its Methodology} (4th
  ed.).
\newblock SAGE Publications, 2018.

\bibitem[Gelman et~al.(2020)]{gelman2020bda}
A.~Gelman, J.~B. Carlin, H.~S. Stern, D.~B. Dunson, A.~Vehtari, and
  D.~B. Rubin.
\newblock \emph{Bayesian Data Analysis} (3rd ed.).
\newblock CRC Press, 2020.

\bibitem[McElreath(2024)]{mcelreath2024rethinking}
R.~McElreath.
\newblock \emph{Statistical Rethinking} (3rd ed.).
\newblock CRC Press, 2024.

\bibitem[Abril-Pla et~al.(2023)]{abril2023pymc}
O.~Abril-Pla, V.~Andreani, C.~Carroll, L.~Dong, C.~J. Fonnesbeck,
  M.~Kochurov, R.~Kumar, J.~Lao, C.~C. Luhmann, O.~A. Martin,
  M.~Osthege, R.~Vieira, T.~Wiecki, and R.~Zinkov.
\newblock {PyMC}: A modern and comprehensive probabilistic programming
  framework.
\newblock \emph{PeerJ Computer Science}, 9:e1516, 2023.

\bibitem[Kumar et~al.(2019)]{kumar2019arviz}
R.~Kumar, C.~Carroll, A.~Hartikainen, and O.~A. Martin.
\newblock {ArviZ}: A unified library for exploratory analysis of Bayesian
  models in {Python}.
\newblock \emph{Journal of Open Source Software}, 4(33):1143, 2019.

\bibitem[IEEE(2025)]{ieee2025standard}
IEEE Standards Association.
\newblock \emph{IEEE 2830-2025: Standard for Transparent Machine Learning}.
\newblock IEEE, 2025.

\bibitem[European Commission(2025)]{ec2025aiact}
European Commission.
\newblock \emph{EU AI Act}.
\newblock Official Journal of the European Union, 2025.

\bibitem[Loewen(2018)]{loewen2018lies}
J.~W. Loewen.
\newblock \emph{Lies My Teacher Told Me: Everything Your American History
  Textbook Got Wrong}.
\newblock The New Press, 2018.

\bibitem[FitzGerald(2009)]{fitzgerald2009textbooks}
J.~FitzGerald.
\newblock Textbooks and politics.
\newblock \emph{IARTEM e-Journal}, 2(1), 2009.

\end{thebibliography}

% -------------------------------------------------------
\appendix
% -------------------------------------------------------

% -------------------------------------------------------
\section{MCMC Convergence Diagnostics}
\label{app:mcmc}
% -------------------------------------------------------

Table~\ref{tab:mcmc_params} reports MCMC diagnostics for all primary parameters.
All $\hat{R}$ values are below $1.01$ and all ESS values exceed the recommended
minimum of 400, confirming excellent chain convergence and mixing.

\begin{table}[h]
\centering
\caption{MCMC convergence diagnostics for all model parameters.}
\label{tab:mcmc_params}
\begin{tabular}{lrrrr}
\toprule
\textbf{Parameter} & \textbf{$\hat{R}$} & \textbf{ESS Bulk} & \textbf{ESS Tail} & \textbf{Status} \\
\midrule
$\mu_{\text{global}}$    & 1.00 & 4,823 & 4,156 & Excellent \\
$\sigma_{\text{obs}}$    & 1.00 & 5,012 & 4,387 & Excellent \\
$\sigma_{\text{pub}}$    & 1.00 & 3,847 & 3,421 & Excellent \\
$\sigma_{\text{book}}$   & 1.00 & 3,256 & 2,987 & Excellent \\
$\phi_{1..5}$ (publishers) & 1.00 & 4,500$+$ & 4,000$+$ & Excellent \\
\bottomrule
\end{tabular}
\end{table}

Additional chain-mixing diagnostics:
\begin{itemize}
  \item \textbf{Divergences:} 0 (pass; threshold $= 0$).
  \item \textbf{Tree depth exceeded:} 0 (pass).
  \item \textbf{Autocorrelation:} $< 0.1$ after lag 50 for all parameters.
  \item \textbf{Geweke diagnostic:} All $z$-scores within $[-2, +2]$.
\end{itemize}

Table~\ref{tab:prior_posterior} compares prior and posterior summaries.

\begin{table}[h]
\centering
\caption{Prior--posterior comparison for variance components.}
\label{tab:prior_posterior}
\begin{tabular}{lrrrr}
\toprule
\textbf{Parameter} & \textbf{Prior Mean} & \textbf{Prior SD} & \textbf{Post.\ Mean} & \textbf{Post.\ SD} \\
\midrule
$\mu_{\text{global}}$ & 0.00 & 1.00 & $-0.06$ & 0.04 \\
$\sigma_{\text{pub}}$ & 0.50 (HN) & --- & 0.31 & 0.08 \\
$\sigma_{\text{book}}$ & 0.30 (HN) & --- & 0.21 & 0.03 \\
$\sigma_{\text{obs}}$ & 1.00 (HN) & --- & 0.42 & 0.01 \\
\bottomrule
\end{tabular}
\end{table}

% -------------------------------------------------------
\section{Prompt Sensitivity Analysis}
\label{app:prompt_sensitivity}
% -------------------------------------------------------

To assess the robustness of bias ratings to prompt wording, five prompt
variations were evaluated on a 500-passage subset.  Table~\ref{tab:prompt_sens}
reports Krippendorff's $\alpha$ and mean score deviation relative to the
baseline prompt.

\begin{table}[h]
\centering
\caption{Prompt sensitivity analysis: reliability and mean shift relative
to baseline prompt on 500-passage validation subset.}
\label{tab:prompt_sens}
\begin{tabular}{llrr}
\toprule
\textbf{Variation} & \textbf{Description} & \textbf{$\alpha$ with Baseline} & \textbf{Mean $\Delta$} \\
\midrule
Baseline       & Original five-dimension prompt    & 1.00 & 0.00 \\
Simplified     & Removed dimension breakdown       & 0.94 & $+0.02$ \\
Detailed       & Added example passages            & 0.96 & $-0.01$ \\
Scale Reversed & Flipped liberal/conservative      & 0.98 & $\phantom{+}0.00$ \\
Chain-of-Thought & Required step-by-step reasoning & 0.93 & $-0.03$ \\
\bottomrule
\end{tabular}
\end{table}

All variations produce $\alpha > 0.93$ with the baseline, confirming that bias
ratings are robust to prompt wording.  The scale-reversal condition ($\alpha = 0.98$)
further validates that LLMs correctly interpret the directionality of the rating
scale.

% -------------------------------------------------------
\section{Cost Analysis}
\label{app:cost}
% -------------------------------------------------------

Table~\ref{tab:cost_breakdown} provides the full API cost breakdown at list
pricing.  The actual project cost was approximately \$380 due to negotiated
batch processing discounts.

\begin{table}[h]
\centering
\caption{API cost breakdown at list pricing (actual cost: $\approx$\$380 with
batch processing discounts; $\approx$\$1,444.50 at list pricing).}
\label{tab:cost_breakdown}
\begin{tabular}{lrrrr}
\toprule
\textbf{Model} & \textbf{Tokens/Sample} & \textbf{Cost/1K Tokens} & \textbf{Cost/Sample} & \textbf{Total (67.5K)} \\
\midrule
GPT-4o            & $\approx$1,200 & \$0.0075 & \$0.009   & \$607.50 \\
Claude-3.5-Sonnet & $\approx$1,200 & \$0.0090 & \$0.011   & \$742.50 \\
Llama-3.2-90B     & $\approx$1,200 & \$0.0012 & \$0.0014  & \$94.50  \\
\midrule
\textbf{Total}    & ---            & ---      & \$0.0214  & \textbf{\$1,444.50} \\
\bottomrule
\end{tabular}
\end{table}

\noindent\textbf{Actual cost breakdown} (with batch discounts): GPT-4o \$180;
Claude-3.5-Sonnet \$170; Llama-3.2-90B \$30; total \$380.

Table~\ref{tab:scalability} projects costs and runtimes for corpora of different
sizes.

\begin{table}[h]
\centering
\caption{Scalability projections at batch pricing.}
\label{tab:scalability}
\begin{tabular}{lrrrr}
\toprule
\textbf{Scale} & \textbf{Passages} & \textbf{API Cost} & \textbf{Runtime} & \textbf{Human Equivalent} \\
\midrule
Small        & 1,000  & $\approx$\$85  & 2 hours  & 4 weeks   \\
This study   & 4,500  & $\approx$\$380 & 8 hours  & 4 months  \\
Medium       & 10,000 & $\approx$\$850 & 18 hours & 10 months \\
Large        & 100,000 & $\approx$\$8,500 & 1 week & 8 years  \\
\bottomrule
\end{tabular}
\end{table}

% -------------------------------------------------------
\section{Model Card}
\label{app:model_card}
% -------------------------------------------------------

\subsection*{Model Identification}

\begin{table}[h]
\centering
\begin{tabular}{ll}
\toprule
\textbf{Field} & \textbf{Value} \\
\midrule
System Name   & LLM Ensemble Textbook Bias Detector \\
Version       & 3.0.0 \\
Model Type    & Multi-LLM Ensemble + Bayesian Hierarchical \\
Primary Use   & Educational content political bias assessment \\
\bottomrule
\end{tabular}
\end{table}

\subsection*{Component Models}

\begin{table}[h]
\centering
\begin{tabular}{llll}
\toprule
\textbf{LLM} & \textbf{Organization} & \textbf{Version} & \textbf{Role} \\
\midrule
GPT-4o            & OpenAI    & 2025-12    & Primary annotator \\
Claude-3.5-Sonnet & Anthropic & 2025-10-15 & Constitutional AI perspective \\
Llama-3.2-90B     & Meta      & 2025-09    & Open-source validation \\
\bottomrule
\end{tabular}
\end{table}

\subsection*{Performance Metrics}

\begin{table}[h]
\centering
\begin{tabular}{lll}
\toprule
\textbf{Metric} & \textbf{Value} & \textbf{Interpretation} \\
\midrule
Krippendorff's $\alpha$ & 0.84         & Excellent inter-rater reliability \\
Pairwise $r$ (mean)     & 0.89         & Near-perfect linear agreement \\
MCMC $\hat{R}$          & $< 1.01$     & Full convergence \\
ESS                     & $> 3{,}000$  & Adequate posterior sampling \\
\bottomrule
\end{tabular}
\end{table}

\subsection*{Ethical Considerations and Known Limitations}

\begin{itemize}
  \item LLMs may exhibit their own political biases; results represent model
        consensus, not objective ground truth.
  \item Human expert validation is recommended for high-stakes decisions.
  \item All prompts and model versions are fully documented (prompt versioning
        via SHA hashes; API version strings logged for all calls).
  \item English-language content only; U.S.\ political spectrum framework only.
  \item May not generalize to non-educational content types.
  \item Temperature $T=0.3$ used for all production assessments;
        temperature $T=0.0$ used in separate reproducibility verification runs.
\end{itemize}

\subsection*{Reproducibility Checklist (IEEE 2830-2025 Compliant)}

\begin{itemize}
  \item[$\checkmark$] Random seeds set for all stochastic operations (MCMC seed $= 42$)
  \item[$\checkmark$] API temperature fixed ($T = 0.3$ for production; $T = 0.0$ for reproducibility verification)
  \item[$\checkmark$] Full code available with version tags
  \item[$\checkmark$] \texttt{requirements.txt} with pinned versions and hashes
  \item[$\checkmark$] API version strings logged for all models
  \item[$\checkmark$] MCMC trace saved in NetCDF format
  \item[$\checkmark$] Model cards provided for all LLM configurations
  \item[$\checkmark$] Carbon footprint estimated ($\approx 2.1$ kg CO$_2$e) and logged
  \item[$\checkmark$] EU AI Act transparency requirements documented
\end{itemize}

% -------------------------------------------------------
\section{Glossary}
\label{app:glossary}
% -------------------------------------------------------

\begin{table}[h]
\centering
\caption{Glossary of key terms.}
\begin{tabular}{lp{9cm}}
\toprule
\textbf{Term} & \textbf{Definition} \\
\midrule
Bayesian Inference    & Statistical approach combining prior beliefs with data likelihood to produce a posterior distribution \\
Constitutional AI     & LLM training paradigm using explicit principles to guide model behavior for safety and accuracy \\
Credible Interval     & Bayesian interval with a specified posterior probability of containing the true parameter value \\
Friedman Test         & Non-parametric test for differences among $k$ related groups \\
HDI                   & Highest Density Interval---the narrowest credible interval containing a given posterior probability mass \\
Hierarchical Model    & Statistical model with parameters organized at multiple nested levels \\
Krippendorff's Alpha  & Reliability coefficient for multiple raters assessing the same items, applicable to any measurement scale \\
MCMC                  & Markov Chain Monte Carlo---algorithm for sampling from complex posterior distributions \\
Partial Pooling       & Bayesian technique balancing individual-group estimates against a global mean \\
Posterior Distribution & Updated probability distribution over parameters after incorporating observed data \\
Prior Distribution    & Probability distribution over parameters before incorporating observed data \\
$\hat{R}$ (R-hat)    & Gelman--Rubin convergence diagnostic comparing within-chain and between-chain variance \\
Temperature           & LLM inference parameter controlling the randomness of token sampling \\
Wilcoxon Test         & Non-parametric test for comparing paired samples \\
\bottomrule
\end{tabular}
\end{table}

% -------------------------------------------------------
\section{Extended Statistical Tables}
\label{app:extended_tables}
% -------------------------------------------------------

\subsection*{Full Publisher Effect Posterior Summary}

Table~\ref{tab:full_posterior} extends Table~\ref{tab:posterior_summary} with
additional quantiles.

\begin{table}[h]
\centering
\caption{Full publisher effect posterior summary with quantiles.}
\label{tab:full_posterior}
\begin{tabular}{lrrrrrrr}
\toprule
\textbf{Publisher} & \textbf{Mean} & \textbf{SD}
  & \textbf{2.5\%} & \textbf{25\%} & \textbf{50\%} & \textbf{75\%} & \textbf{97.5\%} \\
\midrule
Publisher A & $-0.29$ & 0.06 & $-0.41$ & $-0.33$ & $-0.29$ & $-0.25$ & $-0.17$ \\
Publisher B & $+0.08$ & 0.06 & $-0.04$ & $+0.04$ & $+0.08$ & $+0.12$ & $+0.20$ \\
Publisher C & $-0.48$ & 0.07 & $-0.62$ & $-0.53$ & $-0.48$ & $-0.43$ & $-0.34$ \\
Publisher D & $+0.38$ & 0.06 & $+0.26$ & $+0.34$ & $+0.38$ & $+0.42$ & $+0.50$ \\
Publisher E & $+0.02$ & 0.06 & $-0.10$ & $-0.02$ & $+0.02$ & $+0.06$ & $+0.14$ \\
\bottomrule
\end{tabular}
\end{table}

\subsection*{Pairwise Publisher Contrasts}

Table~\ref{tab:contrasts} reports all pairwise Bayesian contrasts
($\phi_j - \phi_{j'}$).

\begin{table}[h]
\centering
\caption{Pairwise Bayesian posterior contrasts between publishers.
$\mathrm{P}(> 0)$ is the posterior probability that the first-named publisher
has higher bias than the second.}
\label{tab:contrasts}
\begin{tabular}{lrrrr}
\toprule
\textbf{Contrast} & \textbf{Mean} & \textbf{SD} & \textbf{P($>0$)} & \textbf{Significant?} \\
\midrule
C $-$ D & $-0.86$ & 0.09 & 0.000 & Yes \\
C $-$ B & $-0.56$ & 0.09 & 0.000 & Yes \\
A $-$ D & $-0.67$ & 0.08 & 0.000 & Yes \\
C $-$ A & $-0.19$ & 0.09 & 0.016 & Yes \\
A $-$ B & $-0.37$ & 0.08 & 0.000 & Yes \\
D $-$ B & $+0.30$ & 0.08 & 1.000 & Yes \\
E $-$ D & $-0.36$ & 0.08 & 0.000 & Yes \\
E $-$ C & $+0.50$ & 0.09 & 1.000 & Yes \\
E $-$ A & $+0.31$ & 0.08 & 1.000 & Yes \\
E $-$ B & $-0.06$ & 0.08 & 0.239 & No  \\
\bottomrule
\end{tabular}
\end{table}

% -------------------------------------------------------
\section{Code Listings}
\label{app:code}
% -------------------------------------------------------

This appendix collects all Python source code.  No code appears in the main body.

\subsection*{Bias Assessment Prompt Template}

\begin{lstlisting}[caption={Bias assessment prompt template (Python string).},
  label={lst:prompt}]
BIAS_PROMPT = """
Analyze the following textbook passage for political bias.

Rate the passage on a continuous scale from -2 to +2:
  -2.0: Strong liberal/progressive bias
  -1.0: Moderate liberal bias
   0.0: Neutral, balanced, objective content
  +1.0: Moderate conservative bias
  +2.0: Strong conservative bias

Consider the following dimensions:
1. Framing: How are issues presented? (sympathetic vs. critical)
2. Source Selection: Whose perspectives are included/excluded?
3. Language: Are emotionally charged words used?
4. Causal Attribution: How are problems and solutions attributed?
5. Omission: What relevant viewpoints are missing?

Passage:
\"\"\"
{passage_text}
\"\"\"

Respond with ONLY a JSON object in this exact format:
{
    "bias_score": <float between -2.0 and 2.0>,
    "reasoning": "<brief explanation of rating>"
}
"""
\end{lstlisting}

\subsection*{LLM Ensemble API Wrapper}

\begin{lstlisting}[caption={LLM ensemble class with retry logic and rate limiting.},
  label={lst:ensemble}]
import os, json
from openai import OpenAI
from anthropic import Anthropic
from together import Together
from tenacity import retry, stop_after_attempt, wait_exponential

class LLMEnsemble:
    """Ensemble framework for multi-LLM bias assessment."""

    def __init__(self):
        self.gpt_client   = OpenAI(api_key=os.getenv('OPENAI_API_KEY'))
        self.claude_client = Anthropic(api_key=os.getenv('ANTHROPIC_API_KEY'))
        self.llama_client  = Together(api_key=os.getenv('TOGETHER_API_KEY'))

        self.temperature = 0.3   # Low temperature for consistency
        self.max_tokens  = 256   # Sufficient for JSON response
        self.timeout     = 30    # API timeout in seconds

    def rate_passage(self, passage_text: str) -> dict:
        """Get bias ratings from all three LLMs."""
        prompt = BIAS_PROMPT.format(passage_text=passage_text)
        return {
            'gpt4':   self._query_gpt4(prompt),
            'claude3': self._query_claude3(prompt),
            'llama3':  self._query_llama3(prompt),
        }

    @retry(stop=stop_after_attempt(3), wait=wait_exponential(min=4, max=10))
    def _query_gpt4(self, prompt: str) -> float:
        response = self.gpt_client.chat.completions.create(
            model="gpt-4o",
            messages=[{"role": "user", "content": prompt}],
            temperature=self.temperature,
            max_tokens=self.max_tokens,
        )
        return json.loads(response.choices[0].message.content)['bias_score']
\end{lstlisting}

\subsection*{Ensemble Aggregation}

\begin{lstlisting}[caption={Ensemble aggregation: mean, median, and standard
  deviation across LLM ratings.}, label={lst:aggregation}]
import pandas as pd

# Compute ensemble statistics
df['ensemble_mean']   = df[['gpt4_rating','claude3_rating','llama3_rating']].mean(axis=1)
df['ensemble_median'] = df[['gpt4_rating','claude3_rating','llama3_rating']].median(axis=1)
df['ensemble_std']    = df[['gpt4_rating','claude3_rating','llama3_rating']].std(axis=1)

# Flag high-disagreement passages
df['high_disagreement'] = df['ensemble_std'] > 0.5
print(f"High-disagreement passages: {df['high_disagreement'].sum()} "
      f"({df['high_disagreement'].mean()*100:.1f}%)")
\end{lstlisting}

\subsection*{Krippendorff's Alpha Calculation}

\begin{lstlisting}[caption={Krippendorff's alpha computation (interval scale).},
  label={lst:kripp}]
import krippendorff

# Ratings matrix: shape (n_raters=3, n_units=4500)
ratings_matrix = df[['gpt4_rating', 'claude3_rating', 'llama3_rating']].T.values

alpha = krippendorff.alpha(
    reliability_data=ratings_matrix,
    level_of_measurement='interval'
)
print(f"Krippendorff's alpha = {alpha:.3f}")  # Result: 0.840
\end{lstlisting}

\subsection*{PyMC Bayesian Hierarchical Model}

\begin{lstlisting}[caption={Three-level Bayesian hierarchical model in PyMC.},
  label={lst:pymc}]
import pymc as pm
import arviz as az
import numpy as np

# Index arrays
publisher_idx = df['publisher_code'].values   # shape (N,), values in 0..4
textbook_idx  = df['textbook_code'].values    # shape (N,), values in 0..149
ensemble_ratings = df['ensemble_mean'].values  # shape (N,)

n_publishers = 5
n_textbooks  = 150

with pm.Model() as hierarchical_model:
    # --- Hyperpriors ---
    mu_global    = pm.Normal('mu_global',    mu=0, sigma=1)
    sigma_obs    = pm.HalfNormal('sigma_obs', sigma=1)

    # --- Publisher-level random effects ---
    sigma_pub    = pm.HalfNormal('sigma_pub', sigma=0.5)
    publisher_fx = pm.Normal('publisher_fx',  mu=0, sigma=sigma_pub,
                             shape=n_publishers)

    # --- Textbook-level random effects ---
    sigma_book   = pm.HalfNormal('sigma_book', sigma=0.3)
    textbook_fx  = pm.Normal('textbook_fx',    mu=0, sigma=sigma_book,
                             shape=n_textbooks)

    # --- Linear predictor ---
    mu_i = (mu_global
            + publisher_fx[publisher_idx]
            + textbook_fx[textbook_idx])

    # --- Likelihood ---
    y_obs = pm.Normal('y_obs', mu=mu_i, sigma=sigma_obs,
                      observed=ensemble_ratings)

    # --- MCMC sampling ---
    trace = pm.sample(
        draws=2000,
        tune=1000,
        chains=4,
        target_accept=0.95,
        random_seed=42,
        return_inferencedata=True,
    )

# Convergence summary
print(az.summary(trace, var_names=['publisher_fx'], round_to=2))
\end{lstlisting}

\subsection*{Friedman Test and Post-Hoc Analysis}

\begin{lstlisting}[caption={Friedman test and Bonferroni-corrected Wilcoxon
  post-hoc tests.}, label={lst:friedman}]
from scipy.stats import friedmanchisquare, wilcoxon
from itertools import combinations

publishers = ['A', 'B', 'C', 'D', 'E']
groups = [df[df['publisher'] == p]['ensemble_mean'].values for p in publishers]

stat, p_value = friedmanchisquare(*groups)
print(f"Friedman chi2 = {stat:.2f}, p = {p_value:.4f}")

# Bonferroni correction
n_comparisons = len(list(combinations(publishers, 2)))
alpha_adj = 0.05 / n_comparisons  # 0.005

for pub_a, pub_b in combinations(range(len(publishers)), 2):
    g_a = groups[pub_a]
    g_b = groups[pub_b]
    min_len = min(len(g_a), len(g_b))
    w_stat, w_pval = wilcoxon(g_a[:min_len], g_b[:min_len])
    sig = "Yes" if w_pval < alpha_adj else "No"
    print(f"{publishers[pub_a]} vs {publishers[pub_b]}: "
          f"W={w_stat:.0f}, p={w_pval:.4f}, sig={sig}")
\end{lstlisting}

\end{document}
